% context: markscheme (official solutions) for 2-5HW_Mental-math
% topic: mental math, no calculator, answers only — metric length prefixes, squares
%        and square roots, (1/2)at^2, m/s and km/h, US units, fractions to decimals,
%        fractions of an hour, a falling stone with g = 10 m/s^2
% convention: answers in bold; "=" joins exact equalities, "\approx" a rounded or
%        approximate value; italic notes give the one-line reasoning and the usual traps
\documentclass[12pt, twoside]{article}
%\documentclass[12pt, twoside]{article}
\usepackage[letterpaper, margin=1in, headsep=0.2in]{geometry}
\setlength{\headheight}{0.6in}
%\usepackage[english]{babel}
\usepackage[utf8]{inputenc}
\usepackage{microtype}
\usepackage{amsmath}
\usepackage{amssymb}
%\usepackage{amsfonts}
\usepackage[nomessages]{fp} %\FPeval{\var-name}{2*sin(pi/6)}
\usepackage{siunitx} %units in math. eg 20\milli\meter
\usepackage{yhmath} % for arcs, overparenth command
\usepackage{tikz} %graphics
\usetikzlibrary{quotes, angles, arrows, arrows.meta}
\usepackage{graphicx} %consider setting \graphicspath{{images/}}
\usepackage{parskip} %no paragraph indent
\usepackage{enumitem}
\usepackage{multicol}
\usepackage{venndiagram}

\usepackage{fancyhdr}
\pagestyle{fancy}
\fancyhf{}
\renewcommand{\headrulewidth}{0pt} % disable the underline of the header
\raggedbottom
\hfuzz=2mm %suppresses overfull box warnings

\usepackage{hyperref}
\usepackage{float}
\setlength{\parskip}{0.2\baselineskip}

\title{Physics}
\author{Chris Huson}
\date{October 2026}

\fancyhead[LE]{\thepage}
\fancyhead[RO]{\thepage}
\fancyhead[LO]{La Scuola d'Italia / Huson / Physics \\* 30 October 2026}

\begin{document}

\subsubsection*{2.5 Homework: Mental Math Solutions}

\emph{Students write only the answer. The notes give the step a student should be
doing in their head, and the mistakes worth looking for.}

\begin{enumerate}[itemsep=0.3cm]

\item Fill in the blanks.
  \begin{multicols}{2}
    \begin{enumerate}[itemsep=0.1cm]
      \item 1.6 km = \textbf{1600 m}
      \item 85 cm = \textbf{0.85 m}
      \item 0.25 m = \textbf{25 cm}
      \item 4 mm = \textbf{0.4 cm}
      \item 2.2 m = \textbf{2200 mm}
      \item 500 mm = \textbf{0.5 m}
    \end{enumerate}
  \end{multicols}
  \emph{Going to a smaller unit makes the number bigger. For (d), 1 cm = 10 mm, so
  4 mm is four tenths of a centimeter.}

\item Squares and square roots.
  \begin{multicols}{4}
    \begin{enumerate}[itemsep=0.1cm]
      \item $3^2 = \mathbf{9}$
      \item $7^2 = \mathbf{49}$
      \item $12^2 = \mathbf{144}$
      \item $20^2 = \mathbf{400}$
      \item $0.5^2 = \mathbf{0.25}$
      \item $0.3^2 = \mathbf{0.09}$
      \item $\sqrt{81} = \mathbf{9}$
      \item $\sqrt{144} = \mathbf{12}$
    \end{enumerate}
  \end{multicols}
  \emph{The trap is (f): $0.3^2$ is $0.3 \times 0.3 = 0.09$, not 0.9. Squaring a number
  less than 1 makes it smaller. Square roots undo squares, which is the last step
  when solving $y = \frac{1}{2}gt^2$ for $t$.}

\item $\frac{1}{2} \times a \times t^2$.
  \begin{multicols}{2}
    \begin{enumerate}[itemsep=0.1cm]
      \item $\frac{1}{2} \times 10 \times 4 = \mathbf{20}$
      \item $\frac{1}{2} \times 10 \times 9 = \mathbf{45}$
      \item $\frac{1}{2} \times 4 \times 25 = \mathbf{50}$
      \item $\frac{1}{2} \times 0.2 \times 100 = \mathbf{10}$
    \end{enumerate}
  \end{multicols}
  \emph{Only $t$ is squared. A student who squares the whole product, or computes
  $(\frac{1}{2} \times 10 \times 2)^2 = 100$ for (a), has the order of operations
  wrong.}

\item Speed units.
  \begin{multicols}{2}
    \begin{enumerate}[itemsep=0.1cm]
      \item 10 m/s = \textbf{36 km/h}
      \item 20 m/s = \textbf{72 km/h}
      \item 108 km/h = \textbf{30 m/s}
      \item 360 km/h = \textbf{100 m/s}
    \end{enumerate}
  \end{multicols}
  \emph{From m/s to km/h, multiply by 3.6; back again, divide. 108 is $3 \times 36$,
  and 360 is $100 \times 3.6$. The number in km/h is always the bigger one.}

\item Circle the best answer.
  \begin{enumerate}[itemsep=0.1cm]
    \item \textbf{(A) 0.3.} \emph{1 ft = 0.3048 m exactly. 3.28 is the number of feet
      in a meter, and 12 is the inches in a foot.}
    \item \textbf{(B) 1.6.} \emph{1 mi $\approx$ 1.609 km. 0.6 is the reverse,
      miles in a kilometer.}
    \item \textbf{(B) 27.} \emph{60 mi/h $\approx$ 97 km/h, which is close to
      100 km/h, and $100 \div 3.6 \approx 28$. The exact value, 26.8 m/s, was given in
      2.3. Choice (D) is the speed in km/h.}
  \end{enumerate}

\newpage

\item Fractions as decimals.
  \begin{multicols}{4}
    \begin{enumerate}[itemsep=0.15cm]
      \item $\dfrac{2}{5} = \mathbf{0.4}$
      \item $\dfrac{4}{5} = \mathbf{0.8}$
      \item $\dfrac{3}{8} = \mathbf{0.375}$
      \item $\dfrac{7}{8} = \mathbf{0.875}$
      \item $\dfrac{1}{20} = \mathbf{0.05}$
      \item $\dfrac{1}{25} = \mathbf{0.04}$
      \item $\dfrac{1}{6} = \mathbf{0.1\overline{6}}$
      \item $\dfrac{5}{6} = \mathbf{0.8\overline{3}}$
    \end{enumerate}
  \end{multicols}
  \emph{Each fifth is 0.2 and each eighth is 0.125. $\frac{1}{20}$ is half of
  $\frac{1}{10}$; $\frac{1}{25}$ is like 4 cents in a dollar. $\frac{1}{6}$ is half
  of $\frac{1}{3}$. Accept $0.1666\ldots$ and $0.8333\ldots$, or 0.17 and 0.83
  written with $\approx$.}

\item Fractions of an hour and a day.
  \begin{multicols}{2}
    \begin{enumerate}[itemsep=0.1cm]
      \item 0.25 hour = \textbf{15 min}
      \item 0.3 hour = \textbf{18 min}
      \item 1.5 min = \textbf{90 s}
      \item 150 s = \textbf{2.5 min}
      \item 12 min = \textbf{0.2 h}
      \item 20 min = $\mathbf{\dfrac{1}{3}}$ \textbf{h}
      \item 1 day = \textbf{1440 min}
      \item 1.75 hours = \textbf{1 h 45 min}
    \end{enumerate}
  \end{multicols}
  \emph{Each tenth of an hour is 6 minutes. For (g), $24 \times 60 = 1440$. As in
  2.1, part (h) is the one to watch: 0.75 h is 45 min, not 75 min.}

\item Class time.
  \begin{enumerate}[itemsep=0.1cm]
    \item \textbf{1.5 h} \quad $2 \times 45 = 90$ min, which is $1\frac{1}{2}$ h.
    \item \textbf{54 h} \quad $36 \times 1.5 = 36 + 18 = 54$.
  \end{enumerate}

\item \textbf{About 6 miles.} \quad $10 \times \dfrac{5}{8} = \dfrac{50}{8}
  = 6\dfrac{1}{4}$.

  \emph{Accept 6 to 6.25 mi; the true value is 6.21 mi.}

\item Falling stone with $g \approx 10$ m/s$^2$.
  \begin{enumerate}[itemsep=0.1cm]
    \item \renewcommand{\arraystretch}{1.4}
      \begin{tabular}[t]{| c | c | c | c | c | c |}
        \hline
        $t$ (s) & 1 & 2 & 3 & 4 & 5 \\ \hline
        speed (m/s) & \textbf{10} & \textbf{20} & \textbf{30} & \textbf{40}
          & \textbf{50} \\ \hline
        distance fallen (m) & \textbf{5} & \textbf{20} & \textbf{45} & \textbf{80}
          & \textbf{125} \\ \hline
      \end{tabular}

      \emph{The speed grows in equal steps; the distance grows as $t^2$, so it is
      $5 \times 1, 4, 9, 16, 25$. With $g = 9.80$ m/s$^2$ the distances are 4.90,
      19.6, 44.1 m, as in the classwork.}
    \item \textbf{8 s.} \quad $5t^2 = 320$, so $t^2 = 64$ and $t = 8$.
    \item \textbf{Between 8 and 9 s.} \quad $5t^2 = 381$ gives $t^2 \approx 76$,
      between $8^2 = 64$ and $9^2 = 81$.

      \emph{With $g = 9.80$ m/s$^2$ the time is 8.82 s. Air resistance would make
      the real fall slower; this is Giancoli's King Kong problem.}
  \end{enumerate}

\end{enumerate}
\end{document}
